\documentclass[11pt]{article}
\usepackage[a4paper,margin=2.4cm]{geometry}
\usepackage[T1]{fontenc}
\usepackage{lmodern}
\usepackage{amsmath,amssymb}
\newcommand{\dphi}[1]{\Delta\phi_{#1}}
\newcommand{\half}{\tfrac{1}{2}}
\title{MOS Capacitor Simulation project\\[4pt]\large Corrected equations (as implemented in the code)}
\date{}
\begin{document}
\maketitle

\section{Model and notation}
The structure is gate -- oxide -- silicon (p-type or n-type) -- grounded bulk contact, in one dimension $x$.
\begin{align}
V_t &= \frac{kT}{q}, & \phi &= \frac{V}{V_t}, & E &= -\frac{dV}{dx}
\end{align}
The grid nodes are $x_0\dots x_N$. Node $0$ is the gate, node $s$ is the oxide--semiconductor interface and
node $N$ is the bulk contact. The oxide mesh is uniform; in the semiconductor the spacing grows geometrically:
\begin{align}
\Delta x_i &= x_{i+1}-x_i, & \Delta x_{s+k} &= a\,b^{k}
\quad\Longrightarrow\quad x_{s+k} = t_{ox} + a\,\frac{b^{k}-1}{b-1}
\end{align}
Each node owns a control volume. Its semiconductor part has width
\begin{align}
h_i = \begin{cases} 0 & i<s \quad\text{(oxide)}\\[2pt]
\half\Delta x_s & i=s \quad\text{(interface: half a cell)}\\[2pt]
\half\left(\Delta x_{i-1}+\Delta x_i\right) & s<i<N \end{cases}
\end{align}
Differences of the normalised potential between neighbouring nodes are written
$\dphi{i}=\phi_{i+1}-\phi_i$.

\section{Current densities (Scharfetter--Gummel)}
\begin{align}
J_n &= q\mu_n nE + qD_n\frac{dn}{dx} = kT\mu_n\left(-n\frac{d\phi}{dx}+\frac{dn}{dx}\right)\\
J_p &= q\mu_p pE - qD_p\frac{dp}{dx} = -kT\mu_p\left(p\frac{d\phi}{dx}+\frac{dp}{dx}\right)
\end{align}
The current is a flux between two nodes, so it is defined on the element $i+\half$ between nodes $i$ and
$i+1$. Taking the potential linear inside an element and writing $n=u\,e^{\phi}$ gives
$J_n = kT\mu_n e^{\phi}\,du/dx$. Integrating $J_n e^{-\phi}$ over the element with constant $J_n$:
\begin{align}
J_{n,i+\half} &= \frac{kT\mu_n}{\Delta x_i}\,B(\dphi{i})\left(n_{i+1}-n_i e^{\dphi{i}}\right)
= \frac{kT\mu_n}{\Delta x_i}\Big[B(\dphi{i})\,n_{i+1}-B(-\dphi{i})\,n_i\Big] \label{eq:jn}
\end{align}
with the Bernoulli function
\begin{align}
B(x)=\frac{x}{e^{x}-1}, \qquad B(x)\,e^{x}=B(-x)
\end{align}
For holes take $p=u\,e^{-\phi}$ (note the minus sign), so $J_p=-kT\mu_p e^{-\phi}\,du/dx$ and
\begin{align}
J_{p,i+\half} &= \frac{kT\mu_p}{\Delta x_i}\,B(\dphi{i})\left(p_i-p_{i+1}e^{\dphi{i}}\right)
= \frac{kT\mu_p}{\Delta x_i}\Big[B(\dphi{i})\,p_i-B(-\dphi{i})\,p_{i+1}\Big] \label{eq:jp}
\end{align}

\section{Steady state (equilibrium and DC bias)}
\subsection{Poisson's equation}
\begin{align}
\frac{d}{dx}\left(\varepsilon\frac{dV}{dx}\right) = -\rho, \qquad \rho = q\left(N_D-N_A+p-n\right)
\end{align}
Integrating over the control volume of node $i$ (there is no charge in the oxide, where $h_i=0$):
\begin{align}
F_i \equiv \varepsilon_{i+\half}\frac{V_{i+1}-V_i}{\Delta x_i}-\varepsilon_{i-\half}\frac{V_i-V_{i-1}}{\Delta x_{i-1}}
+q\,h_i\left(N_D-N_A+p_i-n_i\right)+Q_f\,\delta_{i,s}=0 \label{eq:poisson}
\end{align}
where $\varepsilon_{i+\half}=\varepsilon_{ox}$ for $i<s$ and $\varepsilon_{si}$ for $i\ge s$. At the interface
node $i=s$ the two permittivities differ, which enforces continuity of $\varepsilon E$; the fixed oxide
charge $Q_f$ (per unit area) sits at this node only.

Boundary conditions. $V$ is measured from the bulk, so $V_N=0$. The gate and the semiconductor have
different work functions, so a gate voltage $V_G$ puts the gate node at
\begin{align}
V_0=V_G-\phi_{ms},\qquad \phi_{ms}=\phi_m-\phi_s,\qquad
\phi_s=\chi+\frac{E_g}{2q}-V_t\ln\frac{n_{bulk}}{n_{int}}
\end{align}
with $\phi_m$ the gate work function, $\chi$ the electron affinity and $n_{bulk}$ the electron
concentration at the bulk contact. This covers both substrate types: $N_A>N_D$ is p-type
($n_{bulk}<n_{int}$) and $N_D>N_A$ is n-type. The bands are flat ($V=0$ everywhere) at the flat-band voltage
\begin{align}
V_{FB}=\phi_{ms}-\frac{Q_f}{C_{ox}},\qquad C_{ox}=\frac{\varepsilon_{ox}}{t_{ox}}
\end{align}
Equilibrium is $V_G=0$, which is no longer the flat-band condition.

The equation is non-linear because $n$ and $p$ depend exponentially on $V$. One Newton step
$J\,\Delta V=-F$ is taken with the carriers following the potential at fixed quasi-Fermi levels,
$n\propto e^{V/V_t}$, $p\propto e^{-V/V_t}$:
\begin{align}
J_{i,i\pm1} = \frac{\varepsilon_{i\pm\half}}{\Delta x_{i}\ \text{or}\ \Delta x_{i-1}},\qquad
J_{i,i} = -\frac{\varepsilon_{i+\half}}{\Delta x_i}-\frac{\varepsilon_{i-\half}}{\Delta x_{i-1}}-\frac{q\,h_i\,(n_i+p_i)}{V_t}
\end{align}
followed by $V\leftarrow V+\Delta V$, $n\leftarrow n\,e^{\Delta V/V_t}$, $p\leftarrow p\,e^{-\Delta V/V_t}$.
Steps larger than $V_t$ are damped to $\pm V_t\ln\!\left(1+|\Delta V|/V_t\right)$.

\subsection{Carrier continuity}
\begin{align}
\frac{\partial n}{\partial t}=\frac1q\frac{\partial J_n}{\partial x}-R,\qquad
\frac{\partial p}{\partial t}=-\frac1q\frac{\partial J_p}{\partial x}-R
\end{align}
Recombination is Shockley--Read--Hall, with $n_{int}$ the intrinsic carrier concentration:
\begin{align}
R=\frac{np-n_{int}^2}{S},\qquad S=\tau_p\,(n+n_{int})+\tau_n\,(p+n_{int}) \label{eq:srh}
\end{align}
(A relaxation term $(n-n_{eq})/\tau_n$ with a fixed $n_{eq}$ must not be used here: under bias it pulls the
carriers back to their zero-bias values, although the true steady state of a MOS capacitor has $R=0$ with
$np=n_{int}^2$.) In steady state, integrating over the control volume of node $i$:
\begin{align}
\frac1q\left(J_{n,i+\half}-J_{n,i-\half}\right)=h_iR_i,\qquad
-\frac1q\left(J_{p,i+\half}-J_{p,i-\half}\right)=h_iR_i
\end{align}
Inserting \eqref{eq:jn}, \eqref{eq:jp} and dividing by $V_t\mu$, with $S_i$ and the other carrier taken
from the previous iteration:
\begin{align}
\alpha_{i-1}n_{i-1}+\alpha_in_i+\alpha_{i+1}n_{i+1} &= -\frac{h_i\,n_{int}^2}{V_t\mu_nS_i}\\
\alpha_{i+1}=\frac{B(\dphi{i})}{\Delta x_i},\quad
\alpha_{i-1}&=\frac{B(-\dphi{i-1})}{\Delta x_{i-1}},\quad
\alpha_i=-\left[\frac{B(-\dphi{i})}{\Delta x_i}+\frac{B(\dphi{i-1})}{\Delta x_{i-1}}\right]-\frac{h_i\,p_i}{V_t\mu_nS_i}
\end{align}
\begin{align}
\beta_{i-1}p_{i-1}+\beta_ip_i+\beta_{i+1}p_{i+1} &= -\frac{h_i\,n_{int}^2}{V_t\mu_pS_i}\\
\beta_{i+1}=\frac{B(-\dphi{i})}{\Delta x_i},\quad
\beta_{i-1}&=\frac{B(\dphi{i-1})}{\Delta x_{i-1}},\quad
\beta_i=-\left[\frac{B(\dphi{i})}{\Delta x_i}+\frac{B(-\dphi{i-1})}{\Delta x_{i-1}}\right]-\frac{h_i\,n_i}{V_t\mu_pS_i}
\end{align}
The recombination term enters $\alpha_i$, $\beta_i$ and the right-hand side with a \emph{minus} sign.

Boundary conditions:
\begin{itemize}
\item Interface ($i=s$): no current enters the oxide, $J_{s-\half}=0$. The equations above are used with the
$i-1$ terms dropped and $h_s=\half\Delta x_s$. For $\Delta x_s\to0$ this is $n_{s+1}=n_se^{\dphi{s}}$ and
$p_{s+1}e^{\dphi{s}}=p_s$.
\item Bulk contact ($i=N$): $n_N$, $p_N$ keep the charge-neutral equilibrium values, from
$p-n=N_A-N_D$ and $np=n_{int}^2$.
\end{itemize}
Poisson's equation and the two continuity equations are solved one after the other and repeated until
$V$, $n$ and $p$ stop changing. The gate voltage is moved in steps of at most $0.1$~V, each step starting
from the previous solution.

\subsection{Solving the continuity equations accurately}
Both equations have the form
$\mathcal F_i-\mathcal F_{i-1}-\sigma_iy_i=-g_i$ with flux $\mathcal F_i=u_iy_{i+1}-d_iy_i$ and
$u_i,d_i,\sigma_i,g_i\ge0$. The fluxes are up to $10^{17}$ times larger than the recombination terms, so
ordinary Gaussian elimination loses them to round-off. The elimination is therefore arranged so that it
never subtracts: after eliminating nodes $s\dots i$, node $i$ obeys $u_iy_{i+1}-(d_i+e_i)\,y_i=-b_i$ with
\begin{align}
e_s=\sigma_s,\quad b_s=g_s,\qquad
e_{i+1}=\sigma_{i+1}+\frac{u_i\,e_i}{d_i+e_i},\qquad b_{i+1}=g_{i+1}+\frac{d_i\,b_i}{d_i+e_i}
\end{align}
and the back substitution is $y_i=\left(u_iy_{i+1}+b_i\right)/\left(d_i+e_i\right)$.

\section{Small-signal AC analysis}
A small sinusoidal signal is added to the DC gate bias, and every quantity is written as its DC value plus a
phasor:
\begin{align}
V_G&=V_{DC}+\tilde v\,e^{j\omega t}, & V&=V_0+\tilde V e^{j\omega t}, &
n&=n_0+\tilde n\,e^{j\omega t}, & p&=p_0+\tilde p\,e^{j\omega t}
\end{align}
with $\Delta\tilde\phi_i=(\tilde V_{i+1}-\tilde V_i)/V_t$. Only terms linear in the perturbation are kept,
using
\begin{align}
B(\dphi{i}+\Delta\tilde\phi_ie^{j\omega t})\simeq B(\dphi{0,i})+B'(\dphi{0,i})\,\Delta\tilde\phi_ie^{j\omega t},
\qquad B'(x)=\frac{B(x)\,\big(1-B(-x)\big)}{x}
\end{align}

\subsection{Linearised currents}
From \eqref{eq:jn}, $J_{n,i+\half}=J_{0,i+\half}+\tilde J_{n,i+\half}e^{j\omega t}$ with
\begin{align}
\tilde J_{n,i+\half}&=\frac{kT\mu_n}{\Delta x_i}\Big[B(\dphi{0,i})\,\tilde n_{i+1}-B(-\dphi{0,i})\,\tilde n_i
+g_{n,i}\,\Delta\tilde\phi_i\Big] \label{eq:jnac}\\
g_{n,i}&=B'(\dphi{0,i})\,n_{0,i+1}+B'(-\dphi{0,i})\,n_{0,i}
\end{align}
This is the same as
$B(\dphi{0,i})\big[\tilde n_{i+1}-e^{\dphi{0,i}}(\tilde n_i+n_{0,i}\Delta\tilde\phi_i)\big]
+B'(\dphi{0,i})\,\Delta\tilde\phi_i\big[n_{0,i+1}-n_{0,i}e^{\dphi{0,i}}\big]$,
because $-\frac{d}{dx}\big[B(x)e^{x}\big]=B'(-x)$. From \eqref{eq:jp}:
\begin{align}
\tilde J_{p,i+\half}&=\frac{kT\mu_p}{\Delta x_i}\Big[B(\dphi{0,i})\,\tilde p_i-B(-\dphi{0,i})\,\tilde p_{i+1}
+g_{p,i}\,\Delta\tilde\phi_i\Big] \label{eq:jpac}\\
g_{p,i}&=B'(\dphi{0,i})\,p_{0,i}+B'(-\dphi{0,i})\,p_{0,i+1}
\end{align}

\subsection{Linearised continuity equations}
Linearising \eqref{eq:srh} around the DC solution ($R_0$, $S_0$ are the DC values):
\begin{align}
\tilde R=r_n\tilde n+r_p\tilde p,\qquad r_n=\frac{p_0-R_0\tau_p}{S_0},\qquad r_p=\frac{n_0-R_0\tau_n}{S_0}
\end{align}
The DC parts cancel, and $\partial/\partial t\to j\omega$. For node $i$:
\begin{align}
\frac1q\left(\tilde J_{n,i+\half}-\tilde J_{n,i-\half}\right)-h_i\left(r_{n,i}+j\omega\right)\tilde n_i-h_ir_{p,i}\,\tilde p_i&=0 \label{eq:nac}\\
-\frac1q\left(\tilde J_{p,i+\half}-\tilde J_{p,i-\half}\right)-h_i\left(r_{p,i}+j\omega\right)\tilde p_i-h_ir_{n,i}\,\tilde n_i&=0 \label{eq:pac}
\end{align}
Written out with \eqref{eq:jnac}, and $\alpha_{i\pm1}$ as in the DC case evaluated at the DC potential:
\begin{align}
\alpha_{i+1}\tilde n_{i+1}+\left[\alpha_i^{flux}-\frac{h_i\,(r_{n,i}+j\omega)}{V_t\mu_n}\right]\tilde n_i+\alpha_{i-1}\tilde n_{i-1}
-\frac{h_i\,r_{p,i}}{V_t\mu_n}\tilde p_i+c_{n,i}&=0\\
\alpha_i^{flux}=-\left[\frac{B(-\dphi{0,i})}{\Delta x_i}+\frac{B(\dphi{0,i-1})}{\Delta x_{i-1}}\right],\qquad
c_{n,i}=\frac{g_{n,i}\,\Delta\tilde\phi_i}{\Delta x_i}&-\frac{g_{n,i-1}\,\Delta\tilde\phi_{i-1}}{\Delta x_{i-1}}
\end{align}
and the same for holes with \eqref{eq:jpac}:
\begin{align}
\beta_{i+1}\tilde p_{i+1}+\left[\beta_i^{flux}-\frac{h_i\,(r_{p,i}+j\omega)}{V_t\mu_p}\right]\tilde p_i+\beta_{i-1}\tilde p_{i-1}
-\frac{h_i\,r_{n,i}}{V_t\mu_p}\tilde n_i+c_{p,i}&=0\\
\beta_i^{flux}=-\left[\frac{B(\dphi{0,i})}{\Delta x_i}+\frac{B(-\dphi{0,i-1})}{\Delta x_{i-1}}\right],\qquad
c_{p,i}=-\frac{g_{p,i}\,\Delta\tilde\phi_i}{\Delta x_i}&+\frac{g_{p,i-1}\,\Delta\tilde\phi_{i-1}}{\Delta x_{i-1}}
\end{align}

\subsection{Linearised Poisson equation}
$c_n$, $c_p$ contain $\tilde V$, and $\tilde V$ depends on $\tilde n$, $\tilde p$ through the charge, so
Poisson's equation \eqref{eq:poisson} is linearised and solved together with \eqref{eq:nac}, \eqref{eq:pac}:
\begin{align}
\varepsilon_{i+\half}\frac{\tilde V_{i+1}-\tilde V_i}{\Delta x_i}-\varepsilon_{i-\half}\frac{\tilde V_i-\tilde V_{i-1}}{\Delta x_{i-1}}
+q\,h_i\left(\tilde p_i-\tilde n_i\right)=0 \label{eq:poissonac}
\end{align}

\subsection{Boundary conditions}
\begin{itemize}
\item Gate: $\tilde V_0=\tilde v$. The oxide has no charge, so $\tilde V$ is linear in it and the oxide acts as
the capacitor $C_{ox}=\varepsilon_{ox}/t_{ox}$ between gate and interface.
\item Interface ($i=s$): $\tilde J_{n,s-\half}=\tilde J_{p,s-\half}=0$; equations \eqref{eq:nac},
\eqref{eq:pac} are used with $h_s=\half\Delta x_s$. For $\Delta x_s\to0$ this is $\tilde J_{s+\half}=0$.
\item Bulk contact ($i=N$): $\tilde V_N=0$, $\tilde n_N=\tilde p_N=0$.
\end{itemize}

\subsection{Solution}
The AC problem is linear, so it is solved once per bias point and frequency, with no iteration.
$\phi_{ms}$ and $Q_f$ are constants, so they do not appear in the AC equations.

Solving directly for $\tilde n$, $\tilde p$ is badly conditioned in strong inversion, because the drift
term $g\,\Delta\tilde\phi$ and the diffusion terms of \eqref{eq:jnac} nearly cancel. The code therefore
uses the AC quasi-Fermi potentials $\tilde W_n$, $\tilde W_p$ as unknowns:
\begin{align}
\tilde n=\frac{n_0}{V_t}\left(\tilde V-\tilde W_n\right),\qquad
\tilde p=\frac{p_0}{V_t}\left(\tilde W_p-\tilde V\right)
\end{align}
No DC current flows in a MOS capacitor, so $B(\dphi{0,i})\,n_{0,i+1}=B(-\dphi{0,i})\,n_{0,i}\equiv\bar n_i$,
and then $g_{n,i}=-\bar n_i$. Likewise $B(\dphi{0,i})\,p_{0,i}=B(-\dphi{0,i})\,p_{0,i+1}\equiv\bar p_i=-g_{p,i}$.
Inserting these in \eqref{eq:jnac}, \eqref{eq:jpac}, the cancellation is done analytically and the currents
become plain conductances:
\begin{align}
\frac{\tilde J_{n,i+\half}}{q}=-\frac{\mu_n\bar n_i}{\Delta x_i}\left(\tilde W_{n,i+1}-\tilde W_{n,i}\right),\qquad
\frac{\tilde J_{p,i+\half}}{q}=-\frac{\mu_p\bar p_i}{\Delta x_i}\left(\tilde W_{p,i+1}-\tilde W_{p,i}\right)
\end{align}
These are the same equations as \eqref{eq:nac}, \eqref{eq:pac}, \eqref{eq:poissonac}, only written in other
variables. With $X_i=(\tilde V_i,\tilde W_{n,i},\tilde W_{p,i})^T$, the three equations of node $i$ are
\begin{align}
\big(U_iX_{i+1}-D_iX_i\big)-\big(U_{i-1}X_i-D_{i-1}X_{i-1}\big)-L_iX_i=0
\end{align}
with $3\times3$ complex matrices: $U_i=D_i=\mathrm{diag}\!\left(\varepsilon_{si},\ \mu_n\bar n_i,\ \mu_p\bar p_i\right)/\Delta x_i$
hold the flux coefficients of element $i+\half$ and $L_i$ the
local terms (charge, recombination, $j\omega$). As in the DC case the fluxes are never subtracted from each
other. After eliminating nodes $s\dots i$, node $i$ obeys $U_iX_{i+1}-(D_i+E_i)X_i=-b_i$ with
\begin{align}
E_s&=L_s+\mathrm{diag}(C_{ox},0,0), & b_s&=(C_{ox}\tilde v,0,0)^T\\
E_{i+1}&=L_{i+1}+E_i\,(D_i+E_i)^{-1}U_i, & b_{i+1}&=D_i\,(D_i+E_i)^{-1}b_i
\end{align}
and back substitution $X_i=(D_i+E_i)^{-1}\left(U_iX_{i+1}+b_i\right)$ starting from $X_N=0$.

Two kinds of sweep are done with this solver: a bias sweep ($C$ and $G/\omega$ versus $V_G$ at a few
frequencies) and a frequency sweep ($C$ and $G/\omega$ versus $f$ at a few gate voltages).

\section{Capacitance}
The gate charge per unit area is $Q=\varepsilon_{ox}E_{ox}$, so its phasor is
$\tilde Q=\varepsilon_{ox}(\tilde V_0-\tilde V_1)/\Delta x_0=C_{ox}(\tilde v-\tilde V_s)$ and
\begin{align}
\tilde C=\frac{\tilde Q}{\tilde v}=C_{ox}\left(1-\frac{\tilde V_s}{\tilde v}\right),\qquad
C=\mathrm{Re}\,\tilde C,\qquad G=-\omega\,\mathrm{Im}\,\tilde C
\end{align}
$C$ is the capacitance and $G$ the conductance per unit area at angular frequency $\omega$. At low
frequency the inversion layer follows the signal and $C$ returns to about $C_{ox}$ in inversion; at high
frequency it cannot, and $C$ stays at its minimum.

\section{Corrections made to \texttt{eq.pdf}}
\begin{enumerate}
\item DC recombination term: the sign was wrong in $\alpha_i$, $\beta_i$ and the right-hand sides (old
eqs.\ 16--23). It now enters with a minus sign, consistent with the AC equations.
\item Recombination model: $(n-n_{eq})/\tau$ replaced by SRH, \eqref{eq:srh}, in both DC and AC.
\item $c_n$ (below old eq.\ 40): two sign errors. The $B\,n_{0,i}e^{\dphi{0,i}}\Delta\tilde\phi_i/\Delta x_i$
term is negative and the $B'$ term of element $i-\half$ is negative.
\item $c_p$: the last bracket was $(p_{0,i}-p_{0,i}e^{\dphi{0,i}})$; it is $(p_{0,i}-p_{0,i+1}e^{\dphi{0,i}})$.
\item Holes: ``take $p=u\,e^{\phi}$'' is $p=u\,e^{-\phi}$.
\item Control volume: $\Delta x_{i-1/2}$ is now defined, as $h_i=\half(\Delta x_{i-1}+\Delta x_i)$.
\item Interface boundary: the half control volume of the interface node is kept, instead of setting
$\tilde J=0$ one element inside the semiconductor.
\item Added: discretised Poisson equation (DC and linearised AC), the coupling of the three AC equations, the
capacitance formula, and the missing factor $q$ in $\rho$.
\end{enumerate}

\end{document}
